\addcontentsline{toc}{chapter}{Tóm tắt nội dung}

\begin{center}
    \Large{\textbf{Tóm tắt nội dung}}
\end{center}
\vspace{1cm}

Đồ án nghiên cứu cách tăng tốc quá trình sinh văn bản của mô hình ngôn ngữ lớn (LLM) trên GPU laptop có dung lượng bộ nhớ hạn chế. Trong giải mã tự hồi quy thông thường, mô hình dự đoán từng token một, khiến thời gian sinh có thể kéo dài. Giải mã suy đoán tìm cách giảm số bước tuần tự bằng cách tạo trước một số token ứng viên rồi dùng mô hình đích kiểm tra các token đó. Tuy nhiên, cách dùng một mô hình draft riêng làm tăng nhu cầu bộ nhớ. Vì vậy, đồ án tập trung khảo sát những phương án tạo ứng viên không cần nạp thêm trọng số mô hình.

Hệ thống hiện thực và so sánh giải mã vanilla, phương pháp bỏ qua một số lớp của chính mô hình đích (layer-skip), Prompt Lookup Decoding (PLD) và bộ định tuyến kết hợp các hành động giải mã. Trong phương pháp layer-skip, các lớp được chọn của mô hình đích được lược bỏ khi tạo draft; trong PLD, draft được tìm từ các chuỗi token đã xuất hiện trong ngữ cảnh. Hệ thống cũng quản lý KV cache để tái sử dụng phần tiền tố đã được xác minh. Cách tổ chức này loại bỏ nhu cầu lưu thêm trọng số draft, nhưng bản thân quá trình tạo draft, kiểm chứng và lưu trạng thái tạm thời vẫn có chi phí tính toán và bộ nhớ.

Đồ án đánh giá hệ thống trên NVIDIA RTX 4050 Laptop GPU với Qwen2.5-3B-Instruct ở định dạng NF4. Benchmark paired gồm 150 prompt thuộc GSM8K, HumanEval và CNN/\allowbreak DailyMail, tối đa 48 token đầu ra cho mỗi prompt.

Trong artifact benchmark đầy đủ, PLD đạt speedup end-to-end trung bình $1.063\times$ (CI bootstrap 95\%: $[1.036,1.092]$), còn router đạt $1.050\times$ ($[1.023,1.079]$) so với vanilla. Tỷ số throughput paired router/PLD là $0.989$ ($[0.985,0.993]$), do đó router chưa vượt PLD trong giao thức này. Exact sequence match của PLD và router là 83.3\%. Kết quả cho thấy PLD có thể tăng throughput trung bình trong cấu hình đã đo, nhưng không bảo đảm đầu ra trùng vanilla hoặc mọi tập dữ liệu đều nhanh hơn. Kết luận chỉ áp dụng cho mô hình, phần cứng và giao thức benchmark nêu trên.

Vì speedup của PLD không bảo đảm đầu ra trùng vanilla, đồ án nghiên cứu thêm cách làm cho giải mã suy đoán trên mô hình NF4 cho đúng từng token và KV của greedy tuần tự bitsandbytes gốc, không dùng thêm trọng số draft. Vòng giải mã B2 gồm các phép attention, chuẩn hóa và LM head gọi từng hàng với đúng hình dạng của bước một-token (chính các toán tử gây lệch khi kiểm tra theo khối), quay về bước $M{=}1$ gốc khi không có đề xuất, chiếu LM head chỉ tới hàng bị loại đầu tiên, và một kernel Triton xử lý ba hàng tùy chọn chỉ để tăng tốc thêm khoảng $4\%$. Trên hai cohort held-out GSM8K (mỗi cohort 50 prompt, dùng đúng một lần), B2 khớp token của bitsandbytes ở $50/50$ prompt, nhanh hơn bitsandbytes eager khoảng $1.12$--$1.15\times$ và dùng thêm $0.07$~MiB bộ nhớ đỉnh. B2 chậm hơn một baseline \texttt{torch.compile} không exact ($0.91\times$ tốc độ của nó trên held-out), và bằng chứng về tính exact chỉ là thực nghiệm hữu hạn trên một GPU; lập luận theo toán tử còn các phần chưa được chứng minh.

Cách làm của B2 chạy từng hàng nên bỏ lợi ích đọc trọng số một lần. Vì thế đồ án đề xuất thêm một bộ kernel 4-bit \emph{bất biến theo số hàng $M$} (LM head, attention, RMSNorm và GEMM NF4 viết bằng CUDA C++): kết quả mỗi hàng không phụ thuộc $M$ nhờ khối hàng đệm cố định, thứ tự rút gọn cố định, không split-K, không atomics, được kiểm từng bit trên hàng chục nghìn hàng mỗi kernel. Trên bộ kernel này, vòng PLD exact chạy trong CUDA graph cho đúng token và KV từng bit của giải mã tuần tự $M{=}1$ cùng kernel trên $1200$ cặp prompt-mô hình (tám mô hình Qwen2.5-0.5B/1.5B/3B, OLMo-2-1B, Phi-3.5-mini, Llama-3.2-1B/3B và Mistral-7B; GSM8K, HumanEval, CNN/DailyMail), nhanh hơn baseline exact $1.14$--$1.32\times$ và nhanh hơn vòng exact chạy từng hàng $1.26$--$1.39\times$. Tham chiếu exact ở đây là bộ kernel mới, không phải bitsandbytes gốc: hai tham chiếu cho cùng chuỗi ở chỉ $30$--$44\%$ số prompt GSM8K, và không phát hiện giảm độ chính xác tác vụ ở cỡ mẫu nhỏ. Một mô hình nháp thần kinh vẫn exact nhưng chậm hơn ($0.66\times$) trên đích 3B này. Tất cả kết quả trên một GPU RTX 4050 (6 GB), mô hình tới 7B.

\vspace{0.5cm}
\noindent\textbf{Từ khóa:} giải mã suy đoán, mô hình ngôn ngữ lớn, KV cache, GPU biên, Prompt Lookup Decoding, lượng tử hóa NF4, giải mã exact, kernel bất biến theo số hàng.
